\honest{Fallacies of Compression}{Eliezer Yudkowsky}{2008}

``The map is not the territory,'' as the saying goes, but you cannot fold the territory into your glove compartment. A map of California at a million to one would need a feature size of 10 nanometres to show two leaves a centimetre apart. Reality is huge and your skull is tiny, so distinct things must share one point on your mental map. From the inside, this feels like seeing one thing.

A young child, or ``a sufficiently ancient Greek philosopher,'' would know only one event called ``sound.'' \nb{Aristotle, in \textsc{On the Soul}, explained sound as a movement of air reaching the ear, and wrote that ``'sound' and 'hearing' are both ambiguous.''} Seeing that there are two events behind one point on the map, I say, is ``a big, difficult scientific challenge.''

Confusing two known things under one label is one mistake. The more dangerous one is having ``no idea whatsoever'' that there are two things: one folder, not two folders with one label. A detective may doubt that Carol really has black hair; it takes a subtler detective to ask whether there are two Carols. ``'Tis easier to question one's facts than one's ontology.'' Making the split feels like splitting an atom, and it often comes with a new word, such as ``Carol-2.''

The obvious modern case is ``consciousness,'' which covers being awake, reflectivity, and the hard problem of subjective experience ``as distinguished by Chalmers.'' Press releases announce that consciousness has been explained because someone studied a 40Hz rhythm. The same confusion drives the bait-and-switch in philosophy: argue about one sense of the word, then apply the conclusion to another. \nb{David Chalmers's 1995 paper, the source of the hard problem, also calls ``consciousness'' ambiguous, lists ``the difference between wakefulness and sleep'' among the things it covers, uses Crick and Koch's oscillations as its example of a theory that explains only the easier phenomena, and calls the result a ``bait-and-switch.'' The post names Chalmers only for the hard problem.}

Expanding the map is, ``I say again,'' a scientific challenge. When something has ``protean and self-contradictory attributes,'' it is ``a good guess'' that your map is cramming several things into one point. Avoiding the word, or coining new ones, often helps to pry them apart.
